% =========================================================================
%  II. Overview of the topic, based on the three references
% =========================================================================
\section{Overview of Antialiasing Oscillators}\label{sec:overview}

V\"alim\"aki and Huovilainen \cite{valimaki2007antialiasing} sort
antialiasing oscillators into three groups.

\subsection{Bandlimited Methods}
Additive synthesis builds the sawtooth as a sum of sine waves, one for
each harmonic, and stops at the Nyquist frequency. This avoids
aliasing, but a bass note at 55 Hz already needs 400 sine waves.
Wavetable synthesis reads precomputed bandlimited periods instead, which
becomes difficult when the pitch varies over a wide range
\cite{nam2010efficient}.

\subsection{Quasi-Bandlimited Methods}
All aliasing comes from the jump. The bandlimited impulse train (BLIT) and
the bandlimited step (BLEP) therefore replace each ideal jump with a
smoothed version whose spectrum is limited, usually read from a table
\cite{valimaki2007antialiasing}. The same article introduces PolyBLEP,
which approximates the difference between the smoothed and the ideal step
with a short polynomial. It changes only the two samples next to each
jump and needs no table.

Nam et al.\ \cite{nam2010efficient} build each pulse of a BLIT from a
low-order fractional delay filter, a short filter that can place an
impulse between two samples, and judge the remaining aliases with models
of human hearing. With a third-order
B-spline filter, the waveforms were perceptually free of aliasing up to a
fundamental frequency of 5.8 kHz.

\subsection{Alias-Suppressing Methods}
The differentiated polynomial waveform (DPW) uses the fact that smoother
waveforms have weaker high harmonics. Its second-order form, DPW-2,
squares the naive sawtooth into a parabola. The parabola has no jump and
its harmonics fall as $1/k^2$, so much less energy folds back. The
difference of neighbouring samples then turns it back into a sawtooth.
V\"alim\"aki et al.\ \cite{valimaki2010dpw} extend this idea to higher
orders and found the fourth-order version perceptually free of aliasing
over the whole range of a grand piano. Nam et al.\
\cite{nam2010efficient} found, however, that by their hearing model low
orders alias audibly already at relatively low pitches.
